% Ngân hàng bài mẫu: Phân thức đại số (lớp 8).
% Quy ước:
%   - Mỗi bài: \begin{problem}[id] ... \end{problem}; lời giải lồng bên trong.
%   - id: duy nhất, không dấu, không đổi, không dùng lại. Dạng <lớp><phân môn>-<chủ đề>-<số>.
%   - Mỗi bước lời giải là 1 đoạn, các bước cách nhau 1 dòng trống; mỗi bước có lý do.
%   - Dòng cuối lời giải: Đáp số: $\boxed{...}$
%   - Cách giải khác: thêm \begin{solution}[cách 2].
%   - Metadata bằng comment ngay sau \begin{problem}: % nguon:, % lop:
\documentclass[12pt,a4paper]{article}
\usepackage[utf8]{vietnam}
\usepackage{amsmath,amssymb}
\newenvironment{problem}[1][]{\par\medskip\noindent\textbf{Bài #1.}\ }{\par}
\newenvironment{solution}[1][]{\par\noindent\textit{Lời giải #1.}\par}{\par}

\begin{document}

% ---------------------------------------------------------------- điều kiện xác định, ba phân thức
\begin{problem}[8DS-PS-032]
% nguon: dummy (tự soạn để mở rộng ngân hàng)
% lop: 8
% kc_du_kien: 8DS-PS-01, 8DS-PTNT-01
% do_kho_du_kien: 1
Tìm điều kiện xác định của mỗi phân thức sau:

a) $\dfrac{5}{2x + 6}$;

b) $\dfrac{x - 1}{x^2 - 4x}$;

c) $\dfrac{3x}{x^2 + 1}$.

\begin{solution}
a) Phân thức xác định khi mẫu thức khác $0$: $2x + 6 \neq 0$.

a) Chuyển vế rồi chia hai vế cho $2$: $2x \neq -6$, tức là $x \neq -3$.

b) Phân thức xác định khi $x^2 - 4x \neq 0$. Đặt nhân tử chung $x$: $x^2 - 4x = x(x - 4)$.

b) Tích $x(x - 4)$ khác $0$ khi và chỉ khi cả hai thừa số đều khác $0$: $x \neq 0$ và $x - 4 \neq 0$, tức là $x \neq 0$ và $x \neq 4$.

c) Phân thức xác định khi $x^2 + 1 \neq 0$. Vì $x^2 \geq 0$ với mọi $x$ nên $x^2 + 1 \geq 1 > 0$, tức là mẫu thức luôn khác $0$.

c) Vậy phân thức xác định với mọi giá trị của $x$.

Đáp số: a) $\boxed{x \neq -3}$; b) $\boxed{x \neq 0 \text{ và } x \neq 4}$; c) $\boxed{\text{xác định với mọi } x}$
\end{solution}
\end{problem}

% ---------------------------------------------------------------- rút gọn phân thức đơn giản
\begin{problem}[8DS-PS-033]
% nguon: dummy (tự soạn để mở rộng ngân hàng)
% lop: 8
% kc_du_kien: 8DS-PS-02, 8DS-PTNT-01, 8DS-PTNT-02
% do_kho_du_kien: 1
Rút gọn các phân thức sau:

a) $\dfrac{14x^2y}{21xy^3}$;

b) $\dfrac{4x - 8}{x^2 - 4}$.

\begin{solution}
a) Điều kiện xác định: mẫu $21xy^3 \neq 0$, tức là $x \neq 0$ và $y \neq 0$.

a) Nhân tử chung của tử và mẫu là $7xy$ (ước chung lớn nhất của $14$ và $21$ là $7$, phần biến chung là $xy$). Chia cả tử và mẫu cho $7xy$ (khác $0$ theo điều kiện), theo tính chất cơ bản của phân thức: $\dfrac{14x^2y}{21xy^3} = \dfrac{2x}{3y^2}$.

b) Đặt nhân tử chung ở tử: $4x - 8 = 4(x - 2)$. Phân tích mẫu theo hằng đẳng thức hiệu hai bình phương: $x^2 - 4 = (x - 2)(x + 2)$.

b) Điều kiện xác định: $(x - 2)(x + 2) \neq 0$, tức là $x \neq 2$ và $x \neq -2$.

b) Chia cả tử và mẫu cho $x - 2$ (khác $0$ do $x \neq 2$): $\dfrac{4(x - 2)}{(x - 2)(x + 2)} = \dfrac{4}{x + 2}$.

Đáp số: a) $\boxed{\dfrac{2x}{3y^2} \text{ với } x \neq 0,\ y \neq 0}$; b) $\boxed{\dfrac{4}{x + 2} \text{ với } x \neq 2,\ x \neq -2}$
\end{solution}
\end{problem}

% ---------------------------------------------------------------- cộng, trừ phân thức cùng mẫu và khác mẫu
\begin{problem}[8DS-PS-034]
% nguon: dummy (tự soạn để mở rộng ngân hàng)
% lop: 8
% kc_du_kien: 8DS-PS-03
% do_kho_du_kien: 2
Thực hiện các phép tính sau:

a) $\dfrac{x + 5}{x - 3} - \dfrac{2x + 2}{x - 3}$;

b) $\dfrac{1}{x} + \dfrac{2}{x + 3}$.

\begin{solution}
a) Điều kiện xác định: $x - 3 \neq 0$, tức là $x \neq 3$.

a) Hai phân thức có cùng mẫu $x - 3$, nên trừ các tử và giữ nguyên mẫu (chú ý dấu trừ trước ngoặc): $\dfrac{x + 5}{x - 3} - \dfrac{2x + 2}{x - 3} = \dfrac{x + 5 - 2x - 2}{x - 3} = \dfrac{3 - x}{x - 3}$.

a) Đổi dấu tử: $3 - x = -(x - 3)$. Chia cả tử và mẫu cho $x - 3$ (khác $0$ do $x \neq 3$): $\dfrac{-(x - 3)}{x - 3} = -1$.

b) Điều kiện xác định: các mẫu khác $0$, tức là $x \neq 0$ và $x \neq -3$.

b) Mẫu thức chung là $x(x + 3)$; nhân tử phụ của hai phân thức lần lượt là $x + 3$ và $x$.

b) Quy đồng mẫu rồi cộng các tử, giữ nguyên mẫu chung:
\[ \dfrac{1}{x} + \dfrac{2}{x + 3} = \dfrac{(x + 3) + 2x}{x(x + 3)} = \dfrac{3x + 3}{x(x + 3)} = \dfrac{3(x + 1)}{x(x + 3)}. \]

Đáp số: a) $\boxed{-1 \text{ với } x \neq 3}$; b) $\boxed{\dfrac{3(x + 1)}{x(x + 3)} \text{ với } x \neq 0,\ x \neq -3}$
\end{solution}
\end{problem}

% ---------------------------------------------------------------- chia phân thức, điều kiện của phân thức chia
\begin{problem}[8DS-PS-035]
% nguon: dummy (tự soạn để mở rộng ngân hàng)
% lop: 8
% kc_du_kien: 8DS-PS-04, 8DS-PTNT-01, 8DS-PTNT-02
% do_kho_du_kien: 2
Thực hiện phép tính $\dfrac{x^2 - 9}{x + 1} : \dfrac{x + 3}{3x + 3}$.

\begin{solution}
Phân tích thành nhân tử: $x^2 - 9 = (x - 3)(x + 3)$ (hiệu hai bình phương) và $3x + 3 = 3(x + 1)$ (đặt nhân tử chung).

Điều kiện xác định: các mẫu $x + 1$ và $3(x + 1)$ khác $0$ nên $x \neq -1$; phân thức chia $\dfrac{x + 3}{3x + 3}$ phải khác $0$ nên $x + 3 \neq 0$. Vậy $x \neq -1$ và $x \neq -3$.

Chia cho một phân thức là nhân với phân thức nghịch đảo của nó:
\[ \dfrac{(x - 3)(x + 3)}{x + 1} : \dfrac{x + 3}{3(x + 1)} = \dfrac{(x - 3)(x + 3)}{x + 1} \cdot \dfrac{3(x + 1)}{x + 3} = \dfrac{3(x - 3)(x + 3)(x + 1)}{(x + 1)(x + 3)}. \]

Chia cả tử và mẫu cho nhân tử chung $(x + 1)(x + 3)$ (khác $0$ theo điều kiện), được $3(x - 3)$.

Đáp số: $\boxed{3(x - 3) \text{ với } x \neq -1,\ x \neq -3}$
\end{solution}
\end{problem}

% ---------------------------------------------------------------- hai phân thức bằng nhau, tìm hai đa thức
\begin{problem}[8DS-PS-036]
% nguon: dummy (tự soạn để mở rộng ngân hàng)
% lop: 8
% kc_du_kien: 8DS-PS-06, 8DS-PTNT-01, 8DS-PTNT-02
% do_kho_du_kien: 2
Tìm đa thức $A$ và đa thức $B$ thoả mãn:

a) $\dfrac{A}{x + 4} = \dfrac{x^2 - 16}{x^2 + 8x + 16}$ (với $x \neq -4$);

b) $\dfrac{x^2 - 5x}{B} = \dfrac{x}{x + 2}$ (với các giá trị của $x$ làm các mẫu khác $0$).

\begin{solution}
a) Theo định nghĩa hai phân thức bằng nhau ($\dfrac{A}{B} = \dfrac{C}{D}$ khi $A \cdot D = B \cdot C$), ta có $A \cdot (x^2 + 8x + 16) = (x + 4)(x^2 - 16)$.

a) Viết $x^2 + 8x + 16 = (x + 4)^2$ (bình phương của một tổng) và $x^2 - 16 = (x - 4)(x + 4)$ (hiệu hai bình phương). Do đó $A \cdot (x + 4)^2 = (x + 4)(x - 4)(x + 4) = (x - 4)(x + 4)^2$.

a) Chia hai vế cho $(x + 4)^2$ (khác $0$ do $x \neq -4$): $A = x - 4$.

a) Thử lại: $\dfrac{x^2 - 16}{(x + 4)^2} = \dfrac{(x - 4)(x + 4)}{(x + 4)^2} = \dfrac{x - 4}{x + 4}$ (đúng).

b) Theo định nghĩa hai phân thức bằng nhau: $(x^2 - 5x)(x + 2) = B \cdot x$.

b) Đặt nhân tử chung: $x^2 - 5x = x(x - 5)$, nên $B \cdot x = x(x - 5)(x + 2)$. Chia hai vế cho $x$ (đa thức khác $0$): $B = (x - 5)(x + 2)$.

b) Nhân hai đa thức: $B = x^2 + 2x - 5x - 10 = x^2 - 3x - 10$.

b) Thử lại: với $x \neq 5$ và $x \neq -2$, $\dfrac{x^2 - 5x}{x^2 - 3x - 10} = \dfrac{x(x - 5)}{(x - 5)(x + 2)} = \dfrac{x}{x + 2}$ (đúng).

Đáp số: a) $\boxed{A = x - 4}$; b) $\boxed{B = x^2 - 3x - 10}$
\end{solution}
\end{problem}

% ---------------------------------------------------------------- rút gọn có đổi dấu
\begin{problem}[8DS-PS-037]
% nguon: dummy (tự soạn để mở rộng ngân hàng)
% lop: 8
% kc_du_kien: 8DS-PS-02, 8DS-PTNT-01, 8DS-PTNT-02
% do_kho_du_kien: 2
Rút gọn các phân thức sau:

a) $\dfrac{x^2 - 5x}{25 - x^2}$;

b) $\dfrac{2a^2 - 2ab}{3b - 3a}$.

\begin{solution}
a) Đặt nhân tử chung ở tử: $x^2 - 5x = x(x - 5)$. Phân tích mẫu theo hằng đẳng thức hiệu hai bình phương: $25 - x^2 = (5 - x)(5 + x)$.

a) Điều kiện xác định: $(5 - x)(5 + x) \neq 0$, tức là $x \neq 5$ và $x \neq -5$.

a) Đổi dấu một thừa số của mẫu: $5 - x = -(x - 5)$, nên mẫu bằng $-(x - 5)(x + 5)$.

a) Chia cả tử và mẫu cho $x - 5$ (khác $0$ do $x \neq 5$): $\dfrac{x(x - 5)}{-(x - 5)(x + 5)} = -\dfrac{x}{x + 5}$.

b) Đặt nhân tử chung ở tử: $2a^2 - 2ab = 2a(a - b)$. Đặt nhân tử chung ở mẫu rồi đổi dấu: $3b - 3a = 3(b - a) = -3(a - b)$.

b) Điều kiện xác định: $3b - 3a \neq 0$, tức là $a \neq b$.

b) Chia cả tử và mẫu cho $a - b$ (khác $0$ do $a \neq b$): $\dfrac{2a(a - b)}{-3(a - b)} = -\dfrac{2a}{3}$.

Đáp số: a) $\boxed{-\dfrac{x}{x + 5} \text{ với } x \neq 5,\ x \neq -5}$; b) $\boxed{-\dfrac{2a}{3} \text{ với } a \neq b}$
\end{solution}
\end{problem}

% ---------------------------------------------------------------- rút gọn rồi tính giá trị, điều kiện xác định
\begin{problem}[8DS-PS-038]
% nguon: dummy (tự soạn để mở rộng ngân hàng)
% lop: 8
% kc_du_kien: 8DS-PS-01, 8DS-PS-02, 8DS-PS-05
% do_kho_du_kien: 2
Cho biểu thức $A = \dfrac{x^2 + 2x + 1}{x^2 - 1}$.

a) Tìm điều kiện xác định và rút gọn $A$.

b) Tính giá trị của $A$ tại $x = 3$ và tại $x = -1$ (nếu có).

\begin{solution}
a) Phân tích mẫu theo hằng đẳng thức hiệu hai bình phương: $x^2 - 1 = (x - 1)(x + 1)$. Điều kiện xác định: mẫu khác $0$, tức là $x \neq 1$ và $x \neq -1$.

a) Phân tích tử theo hằng đẳng thức bình phương của một tổng: $x^2 + 2x + 1 = (x + 1)^2$.

a) Chia cả tử và mẫu cho $x + 1$ (khác $0$ do $x \neq -1$): $A = \dfrac{(x + 1)^2}{(x - 1)(x + 1)} = \dfrac{x + 1}{x - 1}$.

b) Giá trị $x = 3$ thoả mãn điều kiện xác định. Thay vào phân thức đã rút gọn: $A = \dfrac{3 + 1}{3 - 1} = 2$.

b) Giá trị $x = -1$ không thoả mãn điều kiện xác định (mẫu $x^2 - 1$ bằng $0$), nên tại $x = -1$ biểu thức $A$ không xác định và không có giá trị, dù phân thức rút gọn $\dfrac{x + 1}{x - 1}$ có giá trị $0$ tại đó.

Đáp số: a) $\boxed{A = \dfrac{x + 1}{x - 1} \text{ với } x \neq 1,\ x \neq -1}$; b) $\boxed{A = 2 \text{ tại } x = 3;\ A \text{ không xác định tại } x = -1}$
\end{solution}
\end{problem}

% ---------------------------------------------------------------- cộng ba phân thức, đổi dấu mẫu
\begin{problem}[8DS-PS-039]
% nguon: dummy (tự soạn để mở rộng ngân hàng)
% lop: 8
% kc_du_kien: 8DS-PS-03, 8DS-PS-02, 8DS-PTNT-02
% do_kho_du_kien: 3
Thực hiện phép tính $\dfrac{x}{x - 3} + \dfrac{x}{x + 3} + \dfrac{18}{9 - x^2}$.

\begin{solution}
Điều kiện xác định: các mẫu khác $0$, tức là $x \neq 3$ và $x \neq -3$ (vì $9 - x^2 = (3 - x)(3 + x)$).

Đổi dấu phân thức thứ ba và phân tích mẫu theo hằng đẳng thức hiệu hai bình phương: $\dfrac{18}{9 - x^2} = \dfrac{-18}{x^2 - 9} = \dfrac{-18}{(x - 3)(x + 3)}$.

Mẫu thức chung là $(x - 3)(x + 3)$; nhân tử phụ của ba phân thức lần lượt là $x + 3$, $x - 3$ và $1$.

Quy đồng mẫu rồi cộng các tử, giữ nguyên mẫu chung:
\[ \dfrac{x}{x - 3} + \dfrac{x}{x + 3} + \dfrac{-18}{(x - 3)(x + 3)} = \dfrac{x(x + 3) + x(x - 3) - 18}{(x - 3)(x + 3)}. \]

Bỏ ngoặc và thu gọn tử: $x^2 + 3x + x^2 - 3x - 18 = 2x^2 - 18 = 2(x^2 - 9) = 2(x - 3)(x + 3)$.

Chia cả tử và mẫu cho $(x - 3)(x + 3)$ (khác $0$ theo điều kiện): kết quả bằng $2$.

Đáp số: $\boxed{2 \text{ với } x \neq 3,\ x \neq -3}$
\end{solution}
\end{problem}

% ---------------------------------------------------------------- bài toán thực tế: hai vòi nước
\begin{problem}[8DS-PS-040]
% nguon: dummy (tự soạn để mở rộng ngân hàng)
% lop: 8
% kc_du_kien: 8DS-PS-03, 8DS-PS-04, 8DS-PS-05
% do_kho_du_kien: 3
Một bể nước đang cạn. Nếu mở riêng vòi I thì sau $x$ giờ bể đầy; nếu mở riêng vòi II thì sau $x + 3$ giờ bể đầy ($x > 0$). Mỗi vòi chảy đều với lưu lượng không đổi.

a) Viết phân thức biểu thị phần bể mà cả hai vòi cùng chảy được trong $1$ giờ.

b) Viết phân thức biểu thị thời gian để hai vòi cùng chảy đầy bể.

c) Tính thời gian đó khi $x = 3$.

\begin{solution}
a) Vì $x > 0$ nên các mẫu $x$ và $x + 3$ đều khác $0$. Trong $1$ giờ, vòi I chảy được $\dfrac{1}{x}$ bể và vòi II chảy được $\dfrac{1}{x + 3}$ bể.

a) Phần bể cả hai vòi chảy được trong $1$ giờ là tổng hai phân thức. Quy đồng với mẫu thức chung $x(x + 3)$ rồi cộng các tử:
\[ \dfrac{1}{x} + \dfrac{1}{x + 3} = \dfrac{(x + 3) + x}{x(x + 3)} = \dfrac{2x + 3}{x(x + 3)} \text{ (bể)}. \]

b) Thời gian bằng toàn bộ bể (là $1$) chia cho phần bể chảy được trong $1$ giờ: $T = 1 : \dfrac{2x + 3}{x(x + 3)}$. Phép chia thực hiện được vì $2x + 3 > 0$ nên phân thức chia khác $0$.

b) Chia cho một phân thức là nhân với phân thức nghịch đảo của nó: $T = \dfrac{x(x + 3)}{2x + 3}$ (giờ).

c) Giá trị $x = 3$ thoả mãn $x > 0$. Thay vào: $T = \dfrac{3 \cdot (3 + 3)}{2 \cdot 3 + 3} = \dfrac{18}{9} = 2$ (giờ).

c) Kiểm tra: vòi I chảy một mình $3$ giờ đầy bể, vòi II chảy một mình $6$ giờ đầy bể, nên mỗi giờ cả hai vòi chảy được $\dfrac{1}{3} + \dfrac{1}{6} = \dfrac{1}{2}$ bể, tức là đầy bể sau $2$ giờ (đúng).

Đáp số: a) $\boxed{\dfrac{2x + 3}{x(x + 3)} \text{ (bể)}}$; b) $\boxed{T = \dfrac{x(x + 3)}{2x + 3} \text{ (giờ)}}$; c) $\boxed{2 \text{ giờ}}$
\end{solution}
\end{problem}

% ---------------------------------------------------------------- chứng minh đẳng thức hai biến, phép chia
\begin{problem}[8DS-PS-041]
% nguon: dummy (tự soạn để mở rộng ngân hàng)
% lop: 8
% kc_du_kien: 8DS-PS-05, 8DS-PS-03, 8DS-PS-04
% do_kho_du_kien: 3
Chứng minh rằng với $a \neq b$ và $a \neq -b$ thì
\[ \left(\dfrac{a}{a - b} - \dfrac{b}{a + b}\right) : \dfrac{a^2 + b^2}{a^2 - b^2} = 1. \]

\begin{solution}
Với $a \neq b$ và $a \neq -b$, các mẫu $a - b$, $a + b$ và $a^2 - b^2 = (a - b)(a + b)$ đều khác $0$. Phân thức chia $\dfrac{a^2 + b^2}{a^2 - b^2}$ khác $0$, vì nếu $a^2 + b^2 = 0$ thì $a = b = 0$, trái với giả thiết $a \neq b$. Do đó vế trái xác định.

Quy đồng trong ngoặc với mẫu thức chung $(a - b)(a + b)$; nhân tử phụ lần lượt là $a + b$ và $a - b$:
\[ \dfrac{a}{a - b} - \dfrac{b}{a + b} = \dfrac{a(a + b) - b(a - b)}{(a - b)(a + b)}. \]

Bỏ ngoặc và thu gọn tử: $a^2 + ab - ab + b^2 = a^2 + b^2$. Mẫu bằng $(a - b)(a + b) = a^2 - b^2$ (hiệu hai bình phương), nên biểu thức trong ngoặc bằng $\dfrac{a^2 + b^2}{a^2 - b^2}$.

Vế trái là thương của hai phân thức giống nhau và phân thức chia khác $0$. Chia cho một phân thức là nhân với phân thức nghịch đảo của nó:
\[ \dfrac{a^2 + b^2}{a^2 - b^2} : \dfrac{a^2 + b^2}{a^2 - b^2} = \dfrac{a^2 + b^2}{a^2 - b^2} \cdot \dfrac{a^2 - b^2}{a^2 + b^2} = \dfrac{(a^2 + b^2)(a^2 - b^2)}{(a^2 - b^2)(a^2 + b^2)}. \]

Chia cả tử và mẫu cho $(a^2 + b^2)(a^2 - b^2)$ (khác $0$): vế trái bằng $1$, đúng bằng vế phải. Đẳng thức được chứng minh.

Đáp số: $\boxed{\left(\dfrac{a}{a - b} - \dfrac{b}{a + b}\right) : \dfrac{a^2 + b^2}{a^2 - b^2} = 1}$
\end{solution}
\end{problem}

% ---------------------------------------------------------------- giá trị nguyên của phân thức
\begin{problem}[8DS-PS-042]
% nguon: dummy (tự soạn để mở rộng ngân hàng)
% lop: 8
% kc_du_kien: 8DS-PS-01, 8DS-PS-05
% do_kho_du_kien: 3
Cho phân thức $P = \dfrac{2x + 5}{x + 1}$.

a) Tìm điều kiện xác định của $P$.

b) Tìm các giá trị nguyên của $x$ để $P$ có giá trị là số nguyên.

\begin{solution}
a) Phân thức xác định khi mẫu thức khác $0$: $x + 1 \neq 0$, tức là $x \neq -1$.

b) Tách tử thành bội của mẫu cộng phần dư: $2x + 5 = 2(x + 1) + 3$. Do đó $P = \dfrac{2(x + 1) + 3}{x + 1} = 2 + \dfrac{3}{x + 1}$.

b) Với $x$ nguyên và $x \neq -1$, số $2$ là số nguyên nên $P$ nguyên khi và chỉ khi $\dfrac{3}{x + 1}$ nguyên, tức là số nguyên $x + 1$ là ước của $3$: $x + 1 \in \{-3;\ -1;\ 1;\ 3\}$.

b) Suy ra $x \in \{-4;\ -2;\ 0;\ 2\}$; cả bốn giá trị đều khác $-1$ (thoả mãn điều kiện xác định).

b) Thử lại: $x = -4$ cho $P = \dfrac{-3}{-3} = 1$; $x = -2$ cho $P = \dfrac{1}{-1} = -1$; $x = 0$ cho $P = 5$; $x = 2$ cho $P = \dfrac{9}{3} = 3$ (đều là số nguyên).

Đáp số: a) $\boxed{x \neq -1}$; b) $\boxed{x \in \{-4;\ -2;\ 0;\ 2\}}$
\end{solution}
\end{problem}

% ---------------------------------------------------------------- rút gọn bằng nhóm hạng tử và hằng đẳng thức
\begin{problem}[8DS-PS-043]
% nguon: dummy (tự soạn để mở rộng ngân hàng)
% lop: 8
% kc_du_kien: 8DS-PS-02, 8DS-PTNT-03, 8DS-PTNT-04
% do_kho_du_kien: 3
Rút gọn phân thức $\dfrac{x^3 - 4x}{x^3 + 2x^2 - 4x - 8}$.

\begin{solution}
Phân tích tử: đặt nhân tử chung $x$ rồi dùng hằng đẳng thức hiệu hai bình phương: $x^3 - 4x = x(x^2 - 4) = x(x - 2)(x + 2)$.

Phân tích mẫu bằng cách nhóm hạng tử rồi đặt nhân tử chung: $x^3 + 2x^2 - 4x - 8 = x^2(x + 2) - 4(x + 2) = (x + 2)(x^2 - 4)$.

Dùng hiệu hai bình phương một lần nữa: $(x + 2)(x^2 - 4) = (x + 2)(x - 2)(x + 2) = (x - 2)(x + 2)^2$.

Điều kiện xác định: mẫu $(x - 2)(x + 2)^2 \neq 0$, tức là $x \neq 2$ và $x \neq -2$.

Chia cả tử và mẫu cho nhân tử chung $(x - 2)(x + 2)$ (khác $0$ theo điều kiện), theo tính chất cơ bản của phân thức: $\dfrac{x(x - 2)(x + 2)}{(x - 2)(x + 2)^2} = \dfrac{x}{x + 2}$.

Đáp số: $\boxed{\dfrac{x}{x + 2} \text{ với } x \neq 2,\ x \neq -2}$
\end{solution}
\end{problem}

% ---------------------------------------------------------------- bài toán thực tế: chi phí trung bình
\begin{problem}[8DS-PS-044]
% nguon: dummy (tự soạn để mở rộng ngân hàng)
% lop: 8
% kc_du_kien: 8DS-PS-05, 8DS-PS-01, 8DS-PT1-01
% do_kho_du_kien: 3
Một xưởng làm đồ gỗ phải trả chi phí cố định $1\,200$ nghìn đồng mỗi ngày (thuê mặt bằng, máy móc). Ngoài ra, mỗi sản phẩm làm ra tốn thêm $40$ nghìn đồng tiền vật liệu. Mỗi ngày xưởng làm $x$ sản phẩm ($x$ là số nguyên dương).

a) Viết phân thức biểu thị chi phí trung bình của một sản phẩm theo $x$.

b) Tính chi phí trung bình của một sản phẩm khi $x = 100$.

c) Xưởng bán mỗi sản phẩm với giá $70$ nghìn đồng. Mỗi ngày xưởng phải làm bao nhiêu sản phẩm để chi phí trung bình của một sản phẩm đúng bằng giá bán?

\begin{solution}
a) Tổng chi phí trong một ngày là $40x + 1\,200$ (nghìn đồng). Chi phí trung bình của một sản phẩm bằng tổng chi phí chia cho số sản phẩm; vì $x > 0$ nên mẫu khác $0$.

a) Đặt nhân tử chung $40$ ở tử: $\dfrac{40x + 1\,200}{x} = \dfrac{40(x + 30)}{x}$ (nghìn đồng).

b) Giá trị $x = 100$ thoả mãn điều kiện. Thay vào: $\dfrac{40 \cdot (100 + 30)}{100} = \dfrac{5\,200}{100} = 52$ (nghìn đồng).

c) Chi phí trung bình bằng giá bán khi $\dfrac{40x + 1\,200}{x} = 70$. Nhân hai vế với $x \neq 0$: $40x + 1\,200 = 70x$.

c) Chuyển vế: $1\,200 = 30x$, suy ra $x = 40$. Giá trị này là số nguyên dương (thoả mãn).

c) Thử lại: tổng chi phí là $40 \cdot 40 + 1\,200 = 2\,800$ nghìn đồng, chia cho $40$ sản phẩm được $70$ nghìn đồng (đúng).

Đáp số: a) $\boxed{\dfrac{40(x + 30)}{x} \text{ (nghìn đồng)}}$; b) $\boxed{52 \text{ nghìn đồng}}$; c) $\boxed{40 \text{ sản phẩm}}$
\end{solution}
\end{problem}

% ---------------------------------------------------------------- biểu thức nhiều phép tính, tính giá trị, giá trị nguyên
\begin{problem}[8DS-PS-045]
% nguon: dummy (tự soạn để mở rộng ngân hàng)
% lop: 8
% kc_du_kien: 8DS-PS-05, 8DS-PS-01, 8DS-PS-03
% do_kho_du_kien: 4
Cho biểu thức
\[ P = \left(\dfrac{3}{x - 1} - \dfrac{x + 5}{x^2 - 1}\right) : \dfrac{2x - 1}{x + 1}. \]
a) Tìm điều kiện xác định và rút gọn $P$.

b) Tính giá trị của $P$ biết $|x - 1| = 2$.

c) Tìm các giá trị nguyên của $x$ để $P$ có giá trị nguyên.

\begin{solution}
a) Các mẫu thức là $x - 1$, $x^2 - 1 = (x - 1)(x + 1)$ và $x + 1$; chúng khác $0$ khi $x \neq 1$ và $x \neq -1$. Phân thức chia $\dfrac{2x - 1}{x + 1}$ phải khác $0$, nên $2x - 1 \neq 0$, tức là $x \neq \dfrac{1}{2}$. Vậy điều kiện xác định là $x \neq 1$, $x \neq -1$, $x \neq \dfrac{1}{2}$.

a) Quy đồng trong ngoặc với mẫu thức chung $(x - 1)(x + 1)$ rồi trừ các tử: $\dfrac{3(x + 1) - (x + 5)}{(x - 1)(x + 1)} = \dfrac{3x + 3 - x - 5}{(x - 1)(x + 1)} = \dfrac{2x - 2}{(x - 1)(x + 1)}$.

a) Đặt nhân tử chung $2x - 2 = 2(x - 1)$, rồi chia cả tử và mẫu cho $x - 1 \neq 0$: biểu thức trong ngoặc bằng $\dfrac{2}{x + 1}$.

a) Chia cho một phân thức là nhân với phân thức nghịch đảo của nó, rồi chia tử và mẫu cho $x + 1 \neq 0$: $P = \dfrac{2}{x + 1} \cdot \dfrac{x + 1}{2x - 1} = \dfrac{2(x + 1)}{(x + 1)(2x - 1)} = \dfrac{2}{2x - 1}$.

b) Từ $|x - 1| = 2$ suy ra $x - 1 = 2$ hoặc $x - 1 = -2$, tức là $x = 3$ hoặc $x = -1$.

b) Đối chiếu điều kiện xác định: $x = -1$ không thoả mãn nên loại; $x = 3$ thoả mãn.

b) Thay $x = 3$ vào biểu thức đã rút gọn: $P = \dfrac{2}{2 \cdot 3 - 1} = \dfrac{2}{5}$.

c) Với $x$ nguyên, $2x - 1$ là số nguyên lẻ. Phân thức $P = \dfrac{2}{2x - 1}$ có giá trị nguyên khi và chỉ khi $2x - 1$ là ước của $2$, tức là $2x - 1 \in \{-2;\ -1;\ 1;\ 2\}$. Vì $2x - 1$ lẻ nên chỉ còn $2x - 1 = -1$ hoặc $2x - 1 = 1$.

c) Giải: $2x - 1 = -1$ cho $x = 0$; $2x - 1 = 1$ cho $x = 1$.

c) Đối chiếu điều kiện xác định: $x = 1$ bị loại; $x = 0$ thoả mãn. Thử lại với biểu thức ban đầu tại $x = 0$: $\left(\dfrac{3}{-1} - \dfrac{5}{-1}\right) : \dfrac{-1}{1} = 2 : (-1) = -2$, là số nguyên (đúng).

Đáp số: a) $\boxed{P = \dfrac{2}{2x - 1} \text{ với } x \neq 1,\ x \neq -1,\ x \neq \dfrac{1}{2}}$; b) $\boxed{P = \dfrac{2}{5}}$; c) $\boxed{x = 0}$
\end{solution}
\end{problem}

% ---------------------------------------------------------------- giá trị nhỏ nhất của biểu thức hữu tỉ
\begin{problem}[8DS-PS-046]
% nguon: dummy (tự soạn để mở rộng ngân hàng)
% lop: 8
% kc_du_kien: 8DS-PS-05, 8DS-HDT-06, 8DS-HDT-01
% do_kho_du_kien: 4
Tìm giá trị nhỏ nhất của biểu thức $D = \dfrac{x^2 + 16}{x}$ với $x > 0$.

\begin{solution}
Vì $x > 0$ nên mẫu khác $0$ và $D$ xác định.

Viết $D = x + \dfrac{16}{x}$. Hai hạng tử $x$ và $\dfrac{16}{x}$ bằng nhau khi $x = 4$, khi đó $D = 8$. Ta chứng minh $D \geq 8$ bằng cách xét hiệu $D - 8$.

Quy đồng với mẫu $x$: $D - 8 = \dfrac{x^2 + 16 - 8x}{x} = \dfrac{x^2 - 8x + 16}{x}$.

Theo hằng đẳng thức bình phương của một hiệu, $x^2 - 8x + 16 = (x - 4)^2$. Vậy $D - 8 = \dfrac{(x - 4)^2}{x}$.

Vì $(x - 4)^2 \geq 0$ và $x > 0$ nên $\dfrac{(x - 4)^2}{x} \geq 0$, suy ra $D \geq 8$ với mọi $x > 0$.

Dấu bằng xảy ra khi $(x - 4)^2 = 0$, tức là $x = 4$ (thoả mãn $x > 0$). Thử lại: $D = \dfrac{16 + 16}{4} = 8$.

Đáp số: $\boxed{\min D = 8 \text{ khi } x = 4}$
\end{solution}
\end{problem}

% ---------------------------------------------------------------- tính giá trị biểu thức từ điều kiện x + 1/x
\begin{problem}[8DS-PS-047]
% nguon: dummy (tự soạn để mở rộng ngân hàng)
% lop: 8
% kc_du_kien: 8DS-PS-05, 8DS-PS-02, 8DS-HDT-01
% do_kho_du_kien: 4
Cho số $x \neq 0$ thoả mãn $x + \dfrac{1}{x} = 3$. Tính giá trị của các biểu thức:

a) $A = x^2 + \dfrac{1}{x^2}$;

b) $B = \dfrac{x^2}{x^4 + x^2 + 1}$.

\begin{solution}
a) Bình phương hai vế của đẳng thức $x + \dfrac{1}{x} = 3$. Theo hằng đẳng thức bình phương của một tổng: $\left(x + \dfrac{1}{x}\right)^2 = x^2 + 2 \cdot x \cdot \dfrac{1}{x} + \dfrac{1}{x^2} = x^2 + 2 + \dfrac{1}{x^2}$.

a) Vế phải là $3^2 = 9$, nên $x^2 + 2 + \dfrac{1}{x^2} = 9$. Suy ra $A = x^2 + \dfrac{1}{x^2} = 9 - 2 = 7$.

b) Mẫu $x^4 + x^2 + 1 \geq 1 > 0$ nên $B$ xác định. Chia cả tử và mẫu cho $x^2 \neq 0$ (tính chất cơ bản của phân thức): $B = \dfrac{1}{\dfrac{x^4 + x^2 + 1}{x^2}} = \dfrac{1}{x^2 + 1 + \dfrac{1}{x^2}}$.

b) Thay kết quả câu a: $x^2 + 1 + \dfrac{1}{x^2} = 7 + 1 = 8$, nên $B = \dfrac{1}{8}$.

b) Kiểm tra giá trị $x$ tồn tại: phương trình $x^2 - 3x + 1 = 0$ (nhân hai vế của giả thiết với $x$) có nghiệm $x = \dfrac{3 \pm \sqrt{5}}{2}$, đều khác $0$.

Đáp số: a) $\boxed{A = 7}$; b) $\boxed{B = \dfrac{1}{8}}$
\end{solution}
\end{problem}

% ---------------------------------------------------------------- tìm hệ số để hai phân thức bằng nhau
\begin{problem}[8DS-PS-048]
% nguon: dummy (tự soạn để mở rộng ngân hàng)
% lop: 8
% kc_du_kien: 8DS-PS-06, 8DS-PTNT-05
% do_kho_du_kien: 4
Tìm các số $a$ và $b$ để với mọi $x \neq 1$ và $x \neq 2$ ta có
\[ \dfrac{x^2 + ax + b}{x^2 - 3x + 2} = \dfrac{x + 3}{x - 2}. \]

\begin{solution}
Phân tích mẫu bằng cách tách $-3x = -x - 2x$ rồi nhóm hạng tử: $x^2 - 3x + 2 = x^2 - x - 2x + 2 = x(x - 1) - 2(x - 1) = (x - 1)(x - 2)$. Với $x \neq 1$ và $x \neq 2$ thì cả hai mẫu $x^2 - 3x + 2$ và $x - 2$ đều khác $0$.

Theo định nghĩa hai phân thức bằng nhau: $(x^2 + ax + b)(x - 2) = (x^2 - 3x + 2)(x + 3) = (x - 1)(x - 2)(x + 3)$.

Chia hai vế cho $x - 2$ (khác $0$ do $x \neq 2$): $x^2 + ax + b = (x - 1)(x + 3) = x^2 + 2x - 3$ với mọi $x \neq 1$, $x \neq 2$.

Hai đa thức bằng nhau tại vô số giá trị của $x$ thì các hệ số của các hạng tử cùng bậc bằng nhau: $a = 2$ và $b = -3$.

Thử lại: với $a = 2$, $b = -3$ ta có $\dfrac{x^2 + 2x - 3}{x^2 - 3x + 2} = \dfrac{(x - 1)(x + 3)}{(x - 1)(x - 2)} = \dfrac{x + 3}{x - 2}$ (đúng).

Đáp số: $\boxed{a = 2,\ b = -3}$
\end{solution}
\end{problem}

% ---------------------------------------------------------------- bài toán thực tế: vận tốc trung bình
\begin{problem}[8DS-PS-049]
% nguon: dummy (tự soạn để mở rộng ngân hàng)
% lop: 8
% kc_du_kien: 8DS-PS-03, 8DS-PS-04, 8DS-PS-05
% do_kho_du_kien: 4
Một người đi xe máy từ A đến B dài $60$ km với vận tốc $x$ (km/h), rồi quay về A với vận tốc nhỏ hơn vận tốc lúc đi $10$ km/h ($x > 10$).

a) Viết phân thức biểu thị tổng thời gian cả đi và về.

b) Vận tốc trung bình của cả hành trình bằng tổng quãng đường chia cho tổng thời gian. Viết phân thức (đã rút gọn) biểu thị vận tốc trung bình theo $x$.

c) Tính vận tốc trung bình khi $x = 30$.

\begin{solution}
a) Vì $x > 10$ nên $x$ và $x - 10$ đều dương, các mẫu khác $0$. Thời gian đi là $\dfrac{60}{x}$ (giờ); vận tốc lúc về là $x - 10$ (km/h) nên thời gian về là $\dfrac{60}{x - 10}$ (giờ).

a) Tổng thời gian là $T = \dfrac{60}{x} + \dfrac{60}{x - 10}$. Quy đồng với mẫu thức chung $x(x - 10)$ rồi cộng các tử:
\[ T = \dfrac{60(x - 10) + 60x}{x(x - 10)} = \dfrac{120x - 600}{x(x - 10)} = \dfrac{120(x - 5)}{x(x - 10)} \text{ (giờ)}. \]

b) Quãng đường cả đi và về là $2 \cdot 60 = 120$ (km). Vận tốc trung bình là $v = 120 : T$; phép chia thực hiện được vì $x > 10$ nên $x - 5 > 0$ và $T \neq 0$.

b) Chia cho một phân thức là nhân với phân thức nghịch đảo của nó, rồi chia tử và mẫu cho $120$: $v = 120 \cdot \dfrac{x(x - 10)}{120(x - 5)} = \dfrac{x(x - 10)}{x - 5}$ (km/h).

c) Giá trị $x = 30$ thoả mãn $x > 10$. Thay vào phân thức đã rút gọn: $v = \dfrac{30 \cdot 20}{25} = 24$ (km/h).

c) Kiểm tra: thời gian đi $\dfrac{60}{30} = 2$ giờ, thời gian về $\dfrac{60}{20} = 3$ giờ, tổng $5$ giờ, và $\dfrac{120}{5} = 24$ (đúng).

Đáp số: a) $\boxed{T = \dfrac{120(x - 5)}{x(x - 10)} \text{ (giờ)}}$; b) $\boxed{v = \dfrac{x(x - 10)}{x - 5} \text{ (km/h)}}$; c) $\boxed{24 \text{ km/h}}$
\end{solution}
\end{problem}

% ---------------------------------------------------------------- điều kiện 1/a + 1/b + 1/c = 0, tính giá trị biểu thức
\begin{problem}[8DS-PS-050]
% nguon: dummy (tự soạn để mở rộng ngân hàng)
% lop: 8
% kc_du_kien: 8DS-PS-05, 8DS-PS-03, 8DS-HDT-03
% do_kho_du_kien: 5
Cho ba số $a$, $b$, $c$ đều khác $0$ và thoả mãn $\dfrac{1}{a} + \dfrac{1}{b} + \dfrac{1}{c} = 0$. Tính giá trị của biểu thức
\[ T = \dfrac{bc}{a^2} + \dfrac{ca}{b^2} + \dfrac{ab}{c^2}. \]

\begin{solution}
Vì $a$, $b$, $c$ đều khác $0$ nên các mẫu của $T$ khác $0$ và $T$ xác định.

Quy đồng vế trái của giả thiết với mẫu thức chung $abc$: $\dfrac{bc + ca + ab}{abc} = 0$. Mẫu $abc \neq 0$ nên tử bằng $0$: $ab + bc + ca = 0$.

Đặt $p = ab$, $q = bc$, $r = ca$. Khi đó $p + q + r = 0$ và $pqr = ab \cdot bc \cdot ca = a^2b^2c^2$.

Quy đồng $T$ với mẫu thức chung $a^2b^2c^2$; nhân tử phụ lần lượt là $b^2c^2$, $c^2a^2$ và $a^2b^2$: $T = \dfrac{b^3c^3 + c^3a^3 + a^3b^3}{a^2b^2c^2} = \dfrac{q^3 + r^3 + p^3}{a^2b^2c^2}$.

Theo hằng đẳng thức lập phương của một tổng, $p^3 + q^3 = (p + q)^3 - 3pq(p + q)$. Vì $p + q = -r$ nên $p^3 + q^3 = -r^3 + 3pqr$, suy ra $p^3 + q^3 + r^3 = 3pqr$.

Vậy $T = \dfrac{3pqr}{a^2b^2c^2} = \dfrac{3a^2b^2c^2}{a^2b^2c^2} = 3$.

Thử lại với $a = 1$, $b = 1$, $c = -\dfrac{1}{2}$: $1 + 1 + (-2) = 0$ và $T = -\dfrac{1}{2} + \left(-\dfrac{1}{2}\right) + 4 = 3$ (đúng).

Đáp số: $\boxed{T = 3}$
\end{solution}
\end{problem}

% ---------------------------------------------------------------- tổng hợp: rút gọn, giải phương trình loại nghiệm, giá trị lớn nhất
\begin{problem}[8DS-PS-051]
% nguon: dummy (tự soạn để mở rộng ngân hàng)
% lop: 8
% kc_du_kien: 8DS-PS-05, 8DS-PS-03, 8DS-HDT-06, 8DS-PTNT-06
% do_kho_du_kien: 5
Cho biểu thức
\[ M = \left(\dfrac{1}{x - 2} - \dfrac{1}{x + 2}\right) \cdot \dfrac{x^2 - 4}{x^2 + 2x + 5}. \]
a) Tìm điều kiện xác định và rút gọn $M$.

b) Tìm $x$ để $M = \dfrac{4}{5}$.

c) Tìm giá trị lớn nhất của $M$.

\begin{solution}
a) Các mẫu thức là $x - 2$, $x + 2$ và $x^2 + 2x + 5$. Theo hằng đẳng thức bình phương của một tổng, $x^2 + 2x + 5 = (x + 1)^2 + 4 \geq 4 > 0$ nên mẫu này luôn khác $0$. Điều kiện xác định: $x \neq 2$ và $x \neq -2$.

a) Quy đồng trong ngoặc với mẫu thức chung $(x - 2)(x + 2) = x^2 - 4$ rồi trừ các tử: $\dfrac{1}{x - 2} - \dfrac{1}{x + 2} = \dfrac{(x + 2) - (x - 2)}{x^2 - 4} = \dfrac{4}{x^2 - 4}$.

a) Nhân hai phân thức rồi chia cả tử và mẫu cho $x^2 - 4 \neq 0$: $M = \dfrac{4}{x^2 - 4} \cdot \dfrac{x^2 - 4}{x^2 + 2x + 5} = \dfrac{4}{x^2 + 2x + 5}$.

b) $M = \dfrac{4}{5}$ khi $\dfrac{4}{x^2 + 2x + 5} = \dfrac{4}{5}$. Hai phân thức có cùng tử $4 \neq 0$ bằng nhau nên mẫu bằng nhau: $x^2 + 2x + 5 = 5$.

b) Chuyển vế và đặt nhân tử chung: $x^2 + 2x = 0$, tức là $x(x + 2) = 0$. Tích bằng $0$ khi một trong hai thừa số bằng $0$: $x = 0$ hoặc $x = -2$.

b) Đối chiếu điều kiện xác định: $x = -2$ không thoả mãn nên loại; $x = 0$ thoả mãn. Thử lại với biểu thức ban đầu tại $x = 0$: $\left(\dfrac{1}{-2} - \dfrac{1}{2}\right) \cdot \dfrac{-4}{5} = (-1) \cdot \dfrac{-4}{5} = \dfrac{4}{5}$ (đúng).

c) Mẫu $x^2 + 2x + 5 = (x + 1)^2 + 4 \geq 4$ với mọi $x$. Phân thức có tử dương $4$ nên mẫu càng lớn thì giá trị càng nhỏ: $M = \dfrac{4}{(x + 1)^2 + 4} \leq \dfrac{4}{4} = 1$.

c) Dấu bằng xảy ra khi $(x + 1)^2 = 0$, tức là $x = -1$. Giá trị này khác $2$ và $-2$ (thoả mãn điều kiện xác định). Thử lại: $M = \dfrac{4}{1 - 2 + 5} = 1$.

Đáp số: a) $\boxed{M = \dfrac{4}{x^2 + 2x + 5} \text{ với } x \neq 2,\ x \neq -2}$; b) $\boxed{x = 0}$; c) $\boxed{\max M = 1 \text{ khi } x = -1}$
\end{solution}
\end{problem}

\end{document}
